\section{Representative Examples}
Given the theoretical framework for downfolding a many orbital (or many-electron) problem to a 
few orbital (or few-electron) problem, we now discuss examples which elucidate the DMD method. 
The examples are as follows:
\begin{itemize}
\item Section~\ref{subsection:3band}: Three-band Hubbard $\rightarrow$ one-band Hubbard at half filling. Demonstrates finding a basis set for the second quantized operators and uses a set of eigenstates directly sampled from the low-energy space to find a one-band model.
\item Section~\ref{subsection:1dhydrogen}: Hydrogen chain $\rightarrow$ one-band Hubbard model at half filling. Demonstrates basis sets for {\it ab initio} systems and the possibility to use this technique to determine the quality of a model for a given physical situation.
\item Section~\ref{subsection:graphene}: Graphene $\rightarrow$ one-band Hubbard model with and without $\sigma$ electrons. Demonstrates using the downfolding procedure to examine the effects of screening due to core electrons. 
\item Section~\ref{subsection:fese}: FeSe molecule $\rightarrow$ $3d,4p,4s$ system. Demonstrates the use of matching pursuit to assess the importance of terms in an effective model and to select compact effective models.
\end{itemize}

In all examples we will highlight the important ingredients associated with DMD. 
First and foremost is the choice of low energy space or energy window i.e. how our database of wave functions was generated. 
Associated with this is the choice of the one body space in terms of which the effective Hamiltonian is expressed. 
Finally, we discuss aspects of the functional forms or parameterizations that are expected to describe our physical problem. 
An important effective Hamiltonian that enters three out of our four representative examples is the one-band or single-band Hubbard model:
\begin{equation}
	H = E_0 -t \;\sum_{\langle i,j \rangle, \eta} \tilde{d}_{i,\eta}^{\dagger} \tilde{d}_{j,\eta} + U \;\sum_{i} \tilde{n}^{i}_{\uparrow} \tilde{n}^{i}_{\downarrow}\,,
\label{eq:oneband}
\end{equation}
where $t$ and $U$ are downfolded (renormalized) parameters, $E_0$ is a constant, $\eta$ is a spin index, 
$\tilde{d}_{i,\eta}$ is the effective one-particle operator associated with spatial orbital (or site) $i$ 
and $n_{i,\eta}=\tilde{d}_{i,\eta}^{\dagger} \tilde{d}_{i,\eta}$ is the corresponding number operator.
$\langle i,j \rangle$ is used to denote nearest neighbor pairs.
%We will sometimes drop the constant energy shift $E_0$ when we write equations like Eq.~\eqref{eq:oneband}.

\subsection{Three-band Hubbard model to one-band Hubbard model at half filling}
\label{subsection:3band} 
Our first example is motivated by the high $T_c$ superconducting cuprates~\cite{Bednorz1986} that 
have parent Mott insulators with rich phase diagrams on electron or hole doping~\cite{Dagotto_RevModPhys, LeeWen_RevModPhys}. 
Many works have been devoted to their model Hamiltonians and corresponding parameter 
values~\cite{tJSpalek, Pavirini, Emery, ZhangRice, Hybertsen_PRB1989, Hybertsen_PRB1990, Kent_Hubbard}. 
A minimal model involving both the copper and oxygen degrees of freedom 
is the three-orbital or three-band Hubbard model, 
\begin{eqnarray}
H &=&    \epsilon_p \sum_{j\in p,\eta} n_{j,\eta} + \epsilon_{d} \sum_{i \in d,\eta}  n_{i,\eta} 
        + t_{pd} \sum_{\langle i\in d ,j \in p \rangle, \eta} \text{sgn}(p_i,d_j) \Big( c_{i,\eta}^{\dagger} c_{j,\eta} + \text{h.c.} \Big) \nonumber \\
          & &   + U_p \sum_{j\in p} n_{j,\uparrow} n_{j,\downarrow} + U_d \sum_{i\in d} n_{i,\uparrow} n_{i,\downarrow} + V_{pd} \sum_{\langle i \in p ,j \in d \rangle} n_j n_i\,,
\end{eqnarray}
where $d_i,p_j$ refer to the  $d_{x^2 - y^2}$ orbitals of copper at site $i$ and $p_x$ or $p_y$ 
oxygen at site $j$, respectively. 
$\text{sgn}(p_i,d_j)$ is the sign of the hopping $t_{pd}$ 
between nearest neighbors, shown schematically in Figure~\ref{fig:threeband}. 
$\epsilon_d$ and $\epsilon_p$ are orbital energies, $U_d$ and $U_p$ are strengths of onsite Hubbard interactions,  
and $V_{pd}$ is the strength of the density-density interactions between a neighboring $p$ and $d$ orbital. 
To simplify, we consider only the case where $\epsilon_p$, $U_d$ and $t_{pd}$ are non zero; $t_{pd}$ is chosen throughout this section to be the typical value of $1.3$ eV to give the reader a sense of overall energy scales. 
Since we work with fixed number of particles we set our reference zero energy to be $\epsilon_d = 0$, thus the charge transfer energy $\Delta \equiv \epsilon_p - \epsilon_d$ equals $\epsilon_p$ in our notation. 
We work in the hole notation; half filling corresponds to two spin-up and two spin-down holes on the $2\times2$ cell.

\begin{figure}[hbt]
\centering
\includegraphics[width=0.9\linewidth]{./Figures/three_band_figure.eps}
\caption
{Schematic for downfolding the three-band Hubbard model on the $2\times2$ cell 
to the one-band Hubbard model. The oxygen orbitals are eliminated to give ``dressed" $d$-like orbitals of the one-band model, 
with modified hopping and interaction parameters. The relationship between the $\tilde{d}$ and the copper and oxygen orbitals is 
encoded by a linear transformation which is parameterized by $\alpha_1$, $\alpha_2$, $\alpha_3$, $\alpha_4$ and $F$ (see Appendix for more details). Here 
the diameter of the circles has been shown to be proportionate to the magnitude of $F$ or $\alpha_i$.}
\label{fig:threeband} 
\end{figure}

It is our objective to determine what one-band Hubbard model~[Eq.~\eqref{eq:oneband}]
``best" describes the three-band data. The effective \textit{d-like} orbitals $\tilde{d}_{i,\eta}$, 
that enter the low energy description are mixtures of copper and oxygen orbitals; this optimal transformation also remains an unknown. 
Thus the model determination involves two aspects (1) what are the composite objects that give a 
compact description of the low energy physics? and (2) given this choice what are the effective interactions between them? 
(A similar problem was posed and solved by one of us in the context of spin systems~\cite{Changlani_percolation}.)
In addition, the best effective Hamiltonian description depends on the energy scale of interest. 
All these issues will be addressed in the remainder of the section. 

We first determine the optimal operators $\tilde{d}_{i,\eta}$, which are constructed from a linear superposition of the $d$ and $p$ orbitals, 
%\HJC{It is not unitary, because we are truncating as well. I am blanking on the precise word} 
\begin{equation}
	\tilde{d}_{i,\eta} = \sum_{j} T_{ij} c_{j,\eta}
\label{eq:dc}
\end{equation}
where $c_{j,\eta}$ is the hole (destruction) operator and refers to either the bare $d$ or $p$ orbitals and $\bf{T}$ is the transformation matrix. 
Further generalizations of this relationship (for example, including higher body terms) are also possible, but have not been considered here. 
For the $2\times2$ unit cell {\bf T} is a $4 \times 12 $ matrix, which we parameterize by 
four distinct parameters. These correspond to mixing of a copper orbital 
with nearest neighbor oxygens ($\alpha_1$), nearest neighbor coppers ($\alpha_2$), next-nearest neighbor oxygens ($\alpha_3$) 
and next-nearest neighbor coppers ($\alpha_4$) as shown schematically in Figure~\ref{fig:threeband}. 
The explicit form of ${\bf T}$ after accounting for the symmetries of the 
lattice has been written out in the Appendix. %These parameters are optimized to minimize a certain cost function, 
%which will be explained shortly. 

%We will now determine the transformation matrix ${\bf T}$. 
All RDMs in the three-band and one-band descriptions are also related via ${\bf T}$; 
the ones that we focus on are evaluated in eigenstate $s$ and are given by,
\begin{subequations}
\begin{eqnarray}
	\langle {\tilde{d}_{i,\eta}}^{\dagger} \tilde{d}_{j,\eta} \rangle_{s} &=& \sum_{mn} T^{*}_{im} \langle {c_{m,\eta}}^{\dagger} c_{n,\eta} \rangle_{s} T_{jn} \label{eq:dmstransformations1} \,,\\
	\langle \tilde{n}_{i,\uparrow} \tilde{n}_{i,\downarrow} \rangle_{s} &=& \sum_{jkmn} T^{*}_{ij} T^{*}_{im} \langle {c_{j,\uparrow}}^{\dagger} {c_{m,\downarrow}}^{\dagger} c_{n,\downarrow} c_{k,\uparrow} \rangle_{s} T_{in} T_{ik}\,.
\label{eq:dmstransformations2}
\end{eqnarray}
\end{subequations}
We determine ${\bf T}$ by demanding two conditions be satisfied, (1) the effective orbitals ($\tilde{d}_{i,\eta}$) 
are orthogonal to each other i.e. $\Big({\bf T} {\bf T}^{\dagger}\Big)_{mn} = \delta_{mn}$
and (2) the sum of all diagonal entries (trace) of the 1-RDM of the effective orbitals for all low energy eigenstates 
equals the number of electrons of a given spin i.e. $\sum_{i} \sum_{\eta} \langle {\tilde{d}_{i,\eta}}^{\dagger} \tilde{d}_{i,\eta} \rangle_{s} = N_{\eta}$. 
These conditions are enforced by minimizing a cost function,
\begin{equation}
C = \sum_{s} \sum_{\eta} \Big( \sum_{i} \langle \tilde{d}_{i,\eta}^{\dagger} \tilde{d}_{i,\eta} \rangle_{s} - N_{\eta} \Big)^{2} + \sum_{mn} ( \Big({\bf T} {\bf T}^{\dagger}\Big)_{mn} -\delta_{mn})^{2}\,.
\label{eq:C}
\end{equation} 
For the $2\times2$ cell, $N_{\uparrow}=N_{\downarrow}=2$ and $i=1,2,3,4$. 
The number of states $s$ was varied from three to six, depending on the energy window of interest. For a selected low energy space, 
we optimize Eq.~\eqref{eq:C} first to determine the optimal one-body operators and then determine the parameters by a 
linear regression fit [Eq.~\eqref{eq:linearfit_descriptor}].

Figure~\ref{fig:varyUdep} shows regimes of the three-band model where 
the lowest six eigenstates are separated from the higher energy manifold; the 
fourth and fifth eigenstates are degenerate. 
In the large $U_d$ limit, charge fluctuations are suppressed and these six 
states correspond to the Hilbert space of $4 \choose 2$ states of the effective spin model in its $S_z=0$ sector.
These states have primarily \textit{d-like} character, an aspect we will verify in this section. 
The eigenstates outside this manifold involve \textit{p-like} excitations which the one-band model is not designed 
to capture. 

We chose the lowest three eigenstates of the three-band model for minimizing 
the cost in Eq.~\eqref{eq:C}. The four dimensional space of parameters of ${\bf T}$ 
was scanned for this purpose. The corresponding trace and orthogonality conditions are simultaneously 
satisfied with only small deviations, confirming the validity of Eq.~\eqref{eq:dc}. 
Importantly, the 1-RDM elements in the transformed basis corresponding to nearest neighbors $\langle \tilde{d}_1^{\dagger} \tilde{d}_2 \rangle_s$ 
already provide estimates for $U/t$ of the effective model. Since the exact knowledge of the corresponding eigenstates of 
the one-band Hubbard model is available for arbitrary $U/t$ by exact diagonalization, we directly look up the $U/t$ with 
the same 1-RDM value. These estimates complement the one obtained by DMD which was carried out 
with the same three low-energy eigenstates, using their energies and 
the computed values of $\langle \tilde{d}_1^{\dagger} \tilde{d}_2 \rangle_s$ 
and $\langle \tilde{n}_{i,\uparrow} \tilde{n}_{i,\downarrow} \rangle_{s}$ from Eqs.~\eqref{eq:dmstransformations1} 
and~\eqref{eq:dmstransformations2}~\footnote{We also mimicked the situation characteristic of \textit{ab initio} 
examples where no eigenstates are generally available. Several non eigenstates were generated as random linear combinations of 
the lowest three eigenstates and input into the DMD procedure, with 
similar outcomes.}. A representative example of our results for $U_{d}/t_{pd}=8$ and $\Delta/t_{pd}=3$ 
has been discussed in the Appendix. 

\renewcommand{\subfigimgone}[3][,]{%
  \setbox1=\hbox{\includegraphics[#1]{#3}}% Store image in box
  \leavevmode\rlap{\usebox1}% Print image
  \rlap{\hspace*{110pt}\vspace*{12pt}\raisebox{\dimexpr\ht1-8\baselineskip}{#2}}% Print label
  \phantom{\usebox1}
}
\renewcommand{\subfigimgtwo}[3][,]{%
  \setbox1=\hbox{\includegraphics[#1]{#3}}% Store image in box
  \leavevmode\rlap{\usebox1}% Print image
  \rlap{\hspace*{95pt}\vspace*{12pt}\raisebox{\dimexpr\ht1-11.3\baselineskip}{#2}}% Print label
  \phantom{\usebox1}
}
\renewcommand{\subfigimgthree}[3][,]{%
  \setbox1=\hbox{\includegraphics[#1]{#3}}% Store image in box
  \leavevmode\rlap{\usebox1}% Print image
  \rlap{\hspace*{125pt}\vspace*{12pt}\raisebox{\dimexpr\ht1-9.7\baselineskip}{#2}}% Print label
  \phantom{\usebox1}
}
\renewcommand{\subfigimgfour}[3][,]{%
  \setbox1=\hbox{\includegraphics[#1]{#3}}% Store image in box
  \leavevmode\rlap{\usebox1}% Print image
  \rlap{\hspace*{95pt}\vspace*{12pt}\raisebox{\dimexpr\ht1-11.5\baselineskip}{#2}}% Print label
  \phantom{\usebox1}
}
\begin{figure}[hbt]
\centering
 \begin{tabular}{@{}p{\linewidth}@{\quad}p{\linewidth}@{}}
\subfigimgone[width=0.47\linewidth]{(A)}{./Figures/spectrum_vs_ep_Ud_8.eps}
\subfigimgtwo[width=0.52\linewidth]{(B)}{./Figures/U_and_hopping_combined_vs_ep_Ud_8.eps}
\subfigimgthree[width=0.47\linewidth]{(C)}{./Figures/spectrum_vs_Ud_ep_3.eps}
\subfigimgfour[width=0.52\linewidth]{(D)}{./Figures/U_and_hopping_combined_vs_Ud_ep_3.eps}
\end{tabular}
\caption{Downfolding of a three-band Hubbard model to an effective one-band Hubbard model of $d$-like orbitals. 
Panels (A) and (C) show the low energy spectra of the three-band model 
(relative to the corresponding ground state) for the cases of (A) fixed $U_d/t_{pd}=8$ and 
varying $\Delta/t_{pd}$ and (C) fixed $\Delta/t_{pd}=3$ and varying $U_d/t_{pd}$. 
In all cases, $t_{pd}$ is set to be $1.3$ eV. 
Panels (B) and (D) show the downfolded parameters for the one-band model corresponding to the three-band parameter 
choices in (A) and (C) respectively.  
The one-band $U/t$ values were obtained either by comparing $\langle \tilde{d}_1^{\dagger} \tilde{d}_2\rangle_s$ 
with the corresponding one-band model eigenstates or by the DMD procedure using 
the lowest three eigenstates. The insets show $t$ obtained from DMD. 
}
\label{fig:varyUdep} 
\end{figure}

Some trends in the one-band description are explored in Figure~\ref{fig:varyUdep} 
by monitoring the downfolded parameters as a function of varying $\Delta/t_{pd}$ and $U_d/t_{pd}$. 
For example, when $U_d/t_{pd}=8$ is fixed and $\Delta/t_{pd}$ is increased, we find that 
the effective hopping $t$ decreases and $U/t$ increases. This is physically reasonable since an increasing difference in the 
single particle energies of the copper and oxygen orbitals makes it energetically unfavorable for holes 
to hop between the two orbitals. When $\Delta/t_{pd}=3$ is fixed and $U_d/t_{pd}$ is increased, $U/t$ increases. 
As one mechanism of avoiding the large $U_d$, the copper orbitals are forced to hybridize more with the oxygen ones; 
on the other hand, hole delocalization is suppressed in a bid to maintain mostly one hole per $\tilde{d}$ due to the larger 
$U/t$. The net result of these effects is that $t$ also increases.%, albeit only marginally; a trend also observed in $\alpha_1$.

An important check for the one-band model is its ability to reproduce the low energy gaps of the three-band model; these have been compared in Figure~\ref{fig:energyfit}. 
For the case of $\Delta/t_{pd}=3$, we observe that for all $U_d/t_{pd}$ the lowest three eigenstates were reproduced well. 
This model also reproduces the states outside of the DMD energy window, although with slightly larger errors. 
Similar trends are seen for the case of $\Delta/t_{pd}=5$, 
with the noticeable difference being that the energy error of the highest state has reduced. 
This also reflects that the parameters obtained from DMD are, in general, dependent on the energy window of interest, a 
point which we will highlight shortly by investigating it systematically. 

\renewcommand{\subfigimgone}[3][,]{%
  \setbox1=\hbox{\includegraphics[#1]{#3}}% Store image in box
  \leavevmode\rlap{\usebox1}% Print image
  \rlap{\hspace*{120pt}\vspace*{1200pt}\raisebox{\dimexpr\ht1-2.0\baselineskip}{#2}}% Print label
  \phantom{\usebox1}
}
\renewcommand{\subfigimgtwo}[3][,]{%
  \setbox1=\hbox{\includegraphics[#1]{#3}}% Store image in box
  \leavevmode\rlap{\usebox1}% Print image
  \rlap{\hspace*{120pt}\vspace*{1200pt}\raisebox{\dimexpr\ht1-2.2\baselineskip}{#2}}% Print label
  \phantom{\usebox1}
}
\begin{figure}[tbh]
\centering
 \begin{tabular}{@{}p{0.90\linewidth}@{\quad}p{\linewidth}@{}}
\subfigimgone[width=0.49\linewidth]{(A)}{./Figures/lowenergy_1and3_vs_Ud_ep_3.eps}
\subfigimgtwo[width=0.49\linewidth]{(B)}{./Figures/lowenergy_1and3_vs_Ud_ep_5.eps}
\end{tabular}
\caption{Comparison of energy spectra of the original three-band model (black circles) and the effective one-band model (red triangles). The three-band model is on the $2\times 2$ cell. 
(A) and (B) show the energy spectra for different parameter sets of the three-band model: (A) $\Delta/t_{pd}=3$ and various $U_{d}/t_{pd}$; (B) $\Delta/t_{pd}=5$ and various $U_{d}/t_{pd}$. 
In all cases, $t_{pd}=1.3$ eV.  }
\label{fig:energyfit} 
\end{figure}	

A promise of downfolding is the reduction of the size of the effective Hilbert space; allowing 
simulations of bigger unit cells to be carried out. To show that this actually works well in practice for the three-band case, 
we consider the $2\sqrt{2} \times 2 \sqrt{2}$ square unit cell, comprising of 8 copper and 16 oxygen orbitals. 
For representative test cases, we performed exact diagonalization calculations at half filling; 
the Hilbert space comprises of 112,911,876 basis states. Roughly 200 Lanczos iterations were carried out, 
enabling convergence of the lowest four energies. We compared the lowest gaps with the 
corresponding calculation on the one-band model on the same square geometry, with a Hilbert space size of only 4,900, 
using the downfolded parameters obtained from the smaller $2 \times 2$ cell. 

Our results are summarized in Figure~\ref{fig:predictivity}. Panel (A) shows six representative parameter sets 
of the three-band model and the corresponding downfolded one-band parameters. Panels (B) and (C) show the lowest three energy 
gaps for representative values of $U_d/t_{pd}=4,8,12$ for $\Delta/t_{pd}=3$ and $\Delta/t_{pd}=5$ respectively. 
In all cases, the agreement between the three-band and one-band models is remarkably good. The 
energy gap error of the lowest gap is within $0.0004$ eV (1\% relative error). The largest error in the third gap is 
of the order of $0.005$ eV (3\% relative error). 
These results indicate the reliability of the downfolding procedure 
and highlight its predictive power. 

\renewcommand{\subfigimgone}[3][,]{%
  \setbox1=\hbox{\includegraphics[#1]{#3}}% Store image in box
  \leavevmode\rlap{\usebox1}% Print image
  \rlap{\hspace*{75pt}\vspace*{100pt}\raisebox{\dimexpr\ht1-10.0\baselineskip}{#2}}% Print label
  \phantom{\usebox1}
}
\renewcommand{\subfigimgtwo}[3][,]{%
  \setbox1=\hbox{\includegraphics[#1]{#3}}% Store image in box
  \leavevmode\rlap{\usebox1}% Print image
  \rlap{\hspace*{75pt}\vspace*{100pt}\raisebox{\dimexpr\ht1-10.0\baselineskip}{#2}}% Print label
  \phantom{\usebox1}
}
\begin{figure}[hbt]
\begin{minipage}{0.42\linewidth}
\centering
 \leavevmode\rlap{\usebox1}% Print image
  \rlap{\hspace*{85pt}\vspace*{100pt}\raisebox{\dimexpr\ht1-5.5\baselineskip}{(A)}}% Print label
  \phantom{\usebox1}
\hbox{% Store image in box
%\begin{tabular}{c|c|c|c||c|c||c|c||c|c}
\quad
\begin{tabular}{c|c||c|c}
\hline
$\Delta/t_{pd}$ & $U_d/t_{pd}$ & $t$ (eV) & $U/t$\\%& $G_{2-1,3}$ & $G_{2-1,1}$ & $G_{3-1,3}$ & $G_{3-1,1}$ & $G_{4-1,3}$ & $G_{4-1,1}$  \\
\hline
\hline
3.0 & 4.0 &  0.2839 & 8.698 \\ %& 0.0465 & 0.0461 & 0.1505 & 0.1489 & 0.2247 & 0.2201  \\ 
3.0 & 8.0 &  0.3025 & 13.45 \\%& 0.0524 & 0.0528 & 0.1682 & 0.1681 & 0.2063 & 0.2020  \\ 
3.0 & 12.0 & 0.3155 & 16.58 \\% & 0.0520 & 0.0524 & 0.1658 & 0.1649 & 0.1920 & 0.1866  \\ 
5.0 & 4.0 &  0.2326 & 15.08 \\%& 0.0395 & 0.0398 & 0.1256 & 0.1258 & 0.1466 & 0.1458  \\ 
5.0 & 8.0 &  0.2501 & 23.75\\% & 0.0344 & 0.0347 & 0.1070 & 0.1070 & 0.1142 & 0.1138  \\ 
5.0 & 12.0 & 0.2647 & 29.89 \\%& 0.0310 & 0.0312 & 0.0956 & 0.0955 & 0.0998 & 0.0993  \\ 
\hline
\end{tabular}}
\end{minipage}
\begin{minipage}{0.57\linewidth}
\subfigimgone[width=0.485\linewidth]{(B)}{./Figures/Predict_8site_Ep_3.eps}
\subfigimgtwo[width=0.485\linewidth]{(C)}{./Figures/Predict_8site_Ep_5.eps}
\end{minipage}
\caption{Multiscale prediction of the effective one-band model. (A) shows the parameters of the effective one-band model obtained from DMD on a small cell, the 4-copper ($2\times2$) cell. 
These parameters were then directly used to predict the energy spectra of a larger cell, the 8-copper ($2\sqrt{2} \times 2\sqrt{2}$) cell. (B) and (C) show the predicted spectra (red triangles) 
in comparison to the exact spectra (black circles) of the three-band model of different parameters (different $\Delta/t_{pd}$ and $U_d/t_{pd}$ where $t_{pd}=1.3$ eV). }
\label{fig:predictivity}
\end{figure} 

Until this point, all our results focused on downfolding using only the lowest three eigenstates of the $2\times2$ cell. 
We now explore the effect of increasing the energy window, by including higher eigenstates, using our 
test example of $U_d/t_{pd}=8$ and $\Delta/t_{pd}=3$. To do so, we now use 
all six low energy eigenstates for optimizing the cost function in Eq.~\eqref{eq:C}. We 
find similar (but not exactly the same) values of $\alpha_i$ compared to 
the case when only the three lowest states were used. The fact that a solution with small cost can be attained 
confirms our expectation that the entire low energy space of six states is consistently 
described by a set of $\tilde{d}_i$ operators. 

\renewcommand{\subfigimgone}[3][,]{%
  \setbox1=\hbox{\includegraphics[#1]{#3}}% Store image in box
  \leavevmode\rlap{\usebox1}% Print image
  \rlap{\hspace*{120pt}\vspace*{12pt}\raisebox{\dimexpr\ht1-10.6\baselineskip}{#2}}% Print label
  \phantom{\usebox1}
}
\renewcommand{\subfigimgtwo}[3][,]{%
  \setbox1=\hbox{\includegraphics[#1]{#3}}% Store image in box
  \leavevmode\rlap{\usebox1}% Print image
  \rlap{\hspace*{120pt}\vspace*{12pt}\raisebox{\dimexpr\ht1-10.5\baselineskip}{#2}}% Print label
  \phantom{\usebox1}
}
\begin{figure}[hbt]
\centering
 \begin{tabular}{@{}p{0.90\linewidth}@{\quad}p{\linewidth}@{}}
\subfigimgone[width=0.49\linewidth]{(A)}{./Figures/downfolded_params_diffwindows_ep_3.eps}
\subfigimgtwo[width=0.50\linewidth]{(B)}{./Figures/lowenergygaps_diffwindows_ep_3.eps}
\end{tabular}
\caption{Dependence of the downfolding parameters on the chosen low energy manifold. (A) shows the variation of downfolded parameters $U/t$ and $t$ (inset) 
with low energy manifolds of different energy windows (corresponding to 3,5 and 6 low energy eigenstates); 
(B) shows the energy spectra of the effective one-band model (red triangles) in comparison to the spectra of the original three-band model (black circles). 
The three-band Hubbard model is on the $2\times 2$ cell with $U_d/t_{pd}=8$, $\Delta/t_{pd}=3$, and $t_{pd}=1.3$ eV.
}
\label{fig:windows} 
\end{figure}	

However, as Figure~\ref{fig:windows}(A) shows, the estimates of $U/t$ and $t$ depend on how many eigenstates 
are used in the DMD procedure. This is because the DMD aims to provide the one-band description 
that best describes \textit{all} states in a given window. If the model is not perfect within a given energy window, 
an energy dependent model is expected, consistent with the renormalization group perspective. For our test example, 
increasing the number of eigenstates from three to six changed $U/t$ from $13.8$ to $9.44$ and $t$ from $0.3045$ to $0.2750$ eV.
\footnote{When three states were used for optimizing the orbitals and for the DMD, we found $U/t\approx 13.45$ and $t=0.3025$ eV. 
This is because $\tilde{d}$ are slightly different in the two cases.} 

The features associated with the energy dependence are further confirmed in Figure~\ref{fig:windows}(B). 
which shows a comparison of energy gaps of the three-band and downfolded one-band model on the $2\times2$ cell. 
When only three states are used, the one-band (nearest neighbor) Hubbard model is insufficient for \textit{accurately}
describing states outside the window. When all six states are used, the DMD tries to minimize the error of the 
largest energy gap at the cost of errors in the smaller energy gaps. 
One could of course choose a different parameterization, say with additional next nearest neighbor $t'$, for which is 
may be possible to reduce this energy dependence significantly and thus have a model that describes the smaller 
and larger energy scales equally well.
 
 
